\documentclass[12pt,a4paper]{article}

% Packages
\usepackage[margin=2.54cm]{geometry}
\usepackage{graphicx}
\usepackage{cite}
\usepackage[hidelinks]{hyperref}
\usepackage{parskip}

% Title information
\title{
    Interim Progress Report\\[0.5cm]
    \large Agentic AI Control of Wireless Networks
    Using an AI-Native Network Controller
    with Safety Guardrails
}

\author{
    Tiernan Bassett\\
    Student Number: 22337085\\
    Supervisor: Merim Dzaferagic\\
    Trinity College Dublin\\
    School of Computer Science and Statistics
}

\date{\today}

\begin{document}

\maketitle
\thispagestyle{empty}

\newpage

\tableofcontents

\newpage

\section{Project Description}

\subsection{Project Overview}
This project will investigate the use of agentic artificial intelligence (AI) in controlling wireless networks. It will connect an agentic workflow built in LangGraph \cite{langgraph} to an AI-Native Network Controller (AI-NNC) \cite{dzaferagic2026ai}. Wireless measurements will be collected and configurations will be issued using the OCUDU platform \cite{ocudu}. The LangGraph workflow will include observation, decision, and execution nodes. There will be three control applications running in parallel within the controller: one for reading measurements, one for issuing configurations, and one for performing rollbacks. The rollback mechanism, inspired by previous post-action verification research \cite{tsvetkov2015post}, will allow the agent to revert to a safe configuration if it deems it necessary. Alongside this post-action safety mechanism, there will be command validation. In the AI-NNC, command validation functions act as guardrails in the control applications before any commands are sent to the network \cite{dzaferagic2026ai}. 

The agent will be able to issue configuration commands and trigger rollbacks using Model Context Protocol (MCP) tools \cite{mcp} available in the applications. Another aspect that will be investigated is how the agent's memory is stored and used. The agent will need to store past measurements, configurations, and experiences to avoid repeating mistakes and reduce the risk of issuing harmful configurations. The memory storage system must also support efficient information retrieval and accommodate the increasing volume of measurements over time \cite{gao2025efficient}. Using this architecture, the application of agentic AI will be evaluated in terms of its safety, efficiency, and network performance.

\subsection{Motivation}
In recent years, large language models (LLMs) have become more capable of reasoning about their tasks \cite{xu2025toward}. Developments such as the ReAct framework, in which LLMs can reason and interact with their environment \cite{yao2023react}, have contributed to the rise of agentic AI \cite{wang2024survey}. Alongside this growth in complexity and reasoning capability, there is now a wide range of applications of agentic AI \cite{wang2024survey}. AI agents have the potential to be used in various sectors, such as law, healthcare, and engineering \cite{wang2024survey}. 


Networking is another sector in which AI can be integrated \cite{polese2023understanding}. AI-native networking is the approach in which AI is a foundational element in networking infrastructure \cite{ferrag20266g}. For example, in wireless networks, radio resource scheduling can be performed by conventional algorithms such as Proportional Fair (PF) or Round-Robin (RR) \cite{mamane2022scheduling}. With AI-native networking, some of these traditional algorithms can be replaced with AI where suitable \cite{mamane2022scheduling}. 


One of the motivations behind AI-native networking is that network complexity and service demand have been increasing due to applications such as cloud computing and digital twins \cite{mamane2022scheduling,ferrag20266g}. A traditional algorithm may not always be able to meet the increase in demand and has limited capabilities to reason about complex scenarios \cite{ferrag20266g}, while agentic AI can reason about many aspects of its environment and past experiences \cite{wang2024survey}. As networks move towards 6G and demand increases, AI-native networking is a promising approach \cite{ferrag20266g}. Traditional network controllers may not have the ability to reason about some situations \cite{ferrag20266g}. For example, a combination of acceptable decisions may lead to network performance degradation \cite{tsvetkov2015post}. An agentic controller may be useful in a dynamic wireless network, with its reasoning capabilities potentially allowing it to outperform traditional algorithms \cite{dzaferagic2026ai}. 


However, agentic AI has limitations such as hallucinations and increased latency \cite{wang2024survey}. Because of these limitations, the application of agentic AI in wireless networks needs to be investigated before deployment.
\subsection{Goal and Objectives}

\subsection{Progress to Date}


\section{Literature Review}

\section{Ethical, Gender and Inclusivity Considerations}

\section{Plan of Work}

\bibliographystyle{IEEEtran}
\bibliography{references}

\end{document}